% ---
% id: protocol-compiler/def:declaration
% kind: def
% title: "Protocol declaration"
% summary: >-
%   A protocol declaration fixes the systems, sector map, marginal slot, both
%   measurements, public maps, raw-key map and a typed constraint list,
%   operators on the measured system being block lists. The general format
%   adds factors, Alice sectors, an ensemble form, a transcript and scoped
%   entries.
% prominence: standard
% status: implemented
% uses: []
% mentions: [def:target, def:exact, thm:band-a, def:homog, def:spec, thm:sifting-and-key, foundations/def:sector-map, rem:block-restriction, req:pass-through, meat-core/def:declaration]
% assumes: []
% refs: []
% provenance: newly-derived
% implements: certiqs_key.protocol_compiler.declaration:CandidateDeclaration
% lean:
% notes: "The completion of the denoted POVM on the kernel of tau_A in eq:ensemblepovm is split among the letters in proportion to tr N_x and not put on the first letter as in transmitter/lem:ensemble-presentation, so that M_x vanishes on a sector exactly when N_x does and key-inertness stays decidable on declared data; part (iv) of that lemma says every such split serves."
% ---

The key-rate engine consumes a protocol declaration. We state it here as this
contract's own definition, with its conditions in
Definitions~\ref{def:exact} and~\ref{def:target}, so that every theorem of
this contract is proved from text inside it. Whether restricting a state
to the \tref[block list]{block lists} of the general format is sound is the subject of
Remark~\ref{rem:block-restriction}. How an entry's form reaches the emitted
declaration is Requirement~\ref{req:pass-through}.

\begin{definition}\label{def:declaration}
A \term[The finite object the compiler emits and the key-rate engine consumes: systems, sector map, marginal slot, measurement families, public maps, raw-key map and constraint list.]{protocol declaration}
is a tuple
\begin{equation}\label{eq:declaration}
  \Decl=\Big(A,\ \tau_A;\ \Qsys,\ \SecSet;\ \{M_x\}_{x\in\Xset};\
  \{F_v\}_{v\in\Vset};\ \Party;\ a_A,\ a_B;\ \phi;\ r,\ \hold;\
  \CList\Big),
\end{equation}
with components as follows. $A$ and $\Qsys$ are finite-dimensional systems,
$\Hs{\Qsys}=\bigoplus_{s\in\SecSet}\Hs{\Qsys_s}$ with a finite index set
$\SecSet$ of sectors, sector projectors $\Pi_s$ and dimensions
$d_s=\dim\Hs{\Qsys_s}$, $d_{\Qsys}=\sum_sd_s$, $d_A=\dim\Hs{A}$. The
\tref{sector map} $\Smap$ is the ordered list
$\big((s,d_s)\big)_{s\in\SecSet}$ of block identifiers with their dimensions,
together with its block kind, a label naming what the identifiers mean
(\xref{foundations/def:sector-map}). It is one object, carried with the
declaration beside the slot $\Qsys,\SecSet$ of \eqref{eq:declaration}, and it
is what fixes $\SecSet$, its total order and $(d_s)_s$.
The
\term[The slot of a declaration that holds Alice's marginal state, or the value marking untrusted-source mode.]{marginal slot}
holds either a state $\tau_A\in\Seq{A}$, in \tref{trusted-source mode}, or
the value $\nomarg$, in \tref{untrusted-source mode}.
$\{M_x\}_{x\in\Xset}$ is a POVM on $A$, given as one matrix per $x$, since $A$
carries no sector decomposition. The words ``state'' and ``POVM'' in the two
sentences above name the intended objects, and they are also exactly the
conditions this contract decides: what a declaration is required to satisfy
is D(i)--D(viii) of Definition~\ref{def:target}, every operator predicate of
which is decided \emph{exactly} on the declared \tref{binary64 numbers},
together with the summation condition of Definition~\ref{def:exact}, without
which a computed sum would not be the sum it is read as.

\emph{Operators on $\Qsys$ are \tref[block list]{block lists}.} An operator
on $\Qsys$ is given as one array $F^{(s)}\in\Lin{\Qsys_s}$ per
$s\in\SecSet$, against the sector map $\Smap$, and that list is the
operator: it is a block list against $\Smap$, the operator it denotes is
$\bigoplus_{s}F^{(s)}$, and the
declaration holds no array larger than a single sector block. Where
$|\SecSet|=1$, say $\SecSet=\{s_0\}$, that one block \emph{is} the whole of
$\Qsys$ and is held and emitted as such, $d_{s_0}=d_{\Qsys}$, as on the
\tref[Band A]{Band-A} route (Theorem~\ref{thm:band-a}\,(i)). Sector block-diagonality is therefore a
property of the \emph{type}: there is no off-sector entry an operator of this
form could carry, and correspondingly nothing to verify. So
$\{F_v\}_{v\in\Vset}$ is a family of such block lists over a finite alphabet,
one list per $v$; $F_v^{(s)}$ denotes its $s$-block. A state $\omega$ on
$A\Qsys$ is likewise a list, one array $\omega^{(s)}\in\Lin{A\Qsys_s}$ per
$s$, and $\omega=\bigoplus_s\omega^{(s)}$ in the same sense. $\Party$ is a
\term[The readout maps of the two parties, with a party label for Alice's system and for each tensor factor of every sector.]{party partition}:
two total maps $\rd{\mathsf A}:\Xset\times\Vset\to\Oset{\mathsf A}$ and
$\rd{\mathsf B}:\Xset\times\Vset\to\Oset{\mathsf B}$ onto declared finite
alphabets, each a
\term[A total map from the outcome pairs to one party's declared finite alphabet, saying what that party reads.]{readout map},
together with a party label for $A$ and, for each $s$, a declared tensor
factorization of $\Hs{\Qsys_s}$ with a party label per factor.
$a_A:\Oset{\mathsf A}\to\CAset$ and $a_B:\Oset{\mathsf B}\to\CBset$, each an
\term[A map from one party's readout alphabet to the alphabet of what that party announces.]{announcement map},
are written jointly
\begin{equation}\label{eq:jointann}
  a(x,v)\coloneqq\big(a_A(\rd{\mathsf A}(x,v)),\,a_B(\rd{\mathsf B}(x,v))\big)
  \in\Kann\coloneqq\CAset\times\CBset ,
\end{equation}
$\phi:\Kann\to\Chatset$ is the acceptance-statistic map,
$r:\Xset\times\Vset\to\Zset$ the raw-key map with
$\Zset=\Zval\sqcup\{\perp\}$ and $\Zval=\{0,\dots,d_Z-1\}$,
$\hold\in\{\mathsf A,\mathsf B\}$ the declared key holder, and
$\CList=\{(\hat O_k,\bowtie_k,\mathrm{rhs}_k)\}_{k=1}^{m}$, $m\ge0$, a finite
constraint list, each entry an operator $\hat O_k$, a relation symbol
$\bowtie_k$ and a right-hand side $\mathrm{rhs}_k$ (which symbols a
well-formed declaration admits is D(vi) of Definition~\ref{def:target}),
with every $\hat O_k$ a block list against the same $\Smap$, Hermitian block
by block.

\emph{The right-hand side of an entry carries a declared tag, and the tag is
part of the datum rather than a property measured of it.} In trusted-source
mode an entry is declared either \tref{homogeneous}, carrying a real scalar
$c_k$, from which the right-hand side it denotes is \emph{constructed} as
$\Theta_k\coloneqq c_k\tau_A$, or \tref{general}, carrying a Hermitian matrix
$\Theta_k$ on $A$. In untrusted-source mode only the scalar form exists and
the right-hand side is a real number $\Theta_k=c_k$. So an entry is written
$(\hat O_k,\bowtie_k,c_k)$ exactly when it is declared homogeneous or is
untrusted-mode, and $(\hat O_k,\bowtie_k,\Theta_k)$ when it is declared
general (Definition~\ref{def:homog}). For these two forms the tag
is the type of the right-hand-side datum, a real scalar or a Hermitian
matrix, which is why Definition~\ref{def:homog} calls them the untagged
forms; the third form, party-local, exists only in the general format below
and is declared by its tag (DX6). We write $\kappa=(c_A,c_B)\in\Kann$, a
value of the joint announcement \eqref{eq:jointann}, for an announcement
class and
\begin{equation}\label{eq:chkey}
  \Chkey\coloneqq\big\{\hat c\in\Chatset:\ \exists\,\kappa\in\phi^{-1}(\hat c)
  \text{ with }\varsigma(\kappa)=\mathrm{pass}\big\}
\end{equation}
for the
\term[Said of an acceptance value that some announcement class passing sifting maps to.]{key-live}
acceptance values. Here $\varsigma$ is the sifting predicate of the
specification's public part (Definition~\ref{def:spec}\,(d)), read as
$\varsigma(\kappa)=\mathrm{fail}$ on every class
$\kappa\notin a(\Xset\times\Vset)$, which no outcome pair realises; on the
realised classes it is, under the hypotheses of
Theorem~\ref{thm:sifting-and-key}, a witness of D(v) of
Definition~\ref{def:target}
(Theorem~\ref{thm:sifting-and-key}\,(ii)). It is \emph{not} a slot of
\eqref{eq:declaration}, so \eqref{eq:chkey} is the set of acceptance values
to which some realised class passing sifting maps, fixed by the
specification and not by the declaration; the key-live set of the consuming
engine, \xref{meat-core/def:declaration}, takes the same convention on the
classes no outcome pair realises.

\emph{The general format.} Everything above is written for the format in
which $A$ carries one sector, one system is measured, Alice's measurement is
declared as the POVM $\{M_x\}$, the public discussion is two messages neither
of which reads the other, and no entry is declared \tref{party-local}; we
call it the
\term[The format with one Alice sector, one measured system, Alice's measurement as a POVM, two messages neither reading the other, and no party-local entry.]{basic format}.
It is the case $|\mathbb{T}|=1$, $|\mathbb{F}|=1$, POVM form, $L=2$ of the
following one, \emph{entry for entry} on every component a derivation reads:
at those values every such object defined in (DX1)--(DX6) is the object of
the same name above, and nothing else is declared. On the party labels of the
tensor factors of each sector block, which no derivation of this module reads
and which the consuming engine's laboratory premise reads as naming the
parties' laboratories, the general format is narrower than the basic one. In general a declaration is the
tuple above with its slots read as follows.
\begin{enumerate}[label=(DX\arabic*),leftmargin=3.2em]
\item \emph{(factorised sector map)} The sector map is an ordered list, each
entry of which is a
\term[One entry of a factorised sector map: a system tag, a party label and a sector map of its own.]{factor}.
Each factor carries a
\term[The label A or Q of a factor, saying whether it is Alice's system or a measured system.]{system tag},
$A$ or $\Qsys$, a party label, and a sector map of its own, an ordered list of
block identifiers with their dimensions; all factors carry one block kind. At
most one factor is tagged $A$, and where there is one it is the first of the
list; at least one factor is tagged $\Qsys$. The factor tagged $A$ fixes a
finite ordered index set $\mathbb{T}$ and a decomposition
$\Hs{A}=\bigoplus_{t\in\mathbb{T}}\Hs{A_t}$ into subspaces spanned by vectors
of the declared basis of $A$, with projectors $\Pi^{A}_t$ and dimensions
$d_{A_t}$, $d_A=\sum_td_{A_t}$; where no factor is tagged $A$, the system
$A$ is still present and is one Alice sector, $\mathbb{T}=\{t_0\}$ and
$A_{t_0}=A$. A factor tagged $\Qsys$ is a
\term[A factor tagged Q: a system one party measures, with its own sector map, alphabet and family.]{measured factor},
and the measured factors form the ordered set $\mathbb{F}$;
factor $f$ fixes a system $\Qsys_f$ with
$\Hs{\Qsys_f}=\bigoplus_{s_f\in\SecSet^{f}}\Hs{\Qsys_{f,s_f}}$, projectors
$\Pi^{f}_{s_f}$ and dimensions $d_{f,s_f}$, and for $s=(s_f)_{f}$
\begin{equation}\label{eq:productmap}
  \Qsys=\bigotimes_{f\in\mathbb{F}}\Qsys_f,\qquad
  \SecSet=\prod_{f\in\mathbb{F}}\SecSet^{f},\qquad
  \Hs{\Qsys_s}=\bigotimes_{f}\Hs{\Qsys_{f,s_f}},\qquad
  \Pi_s=\bigotimes_f\Pi^{f}_{s_f},\qquad d_s=\prod_fd_{f,s_f},
\end{equation}
all products in the order of $\mathbb{F}$, and $\SecSet$ ordered with the
last factor fastest, a product identifier being the factors' identifiers
joined in that order. Each pair $(t,s)\in\mathbb{T}\times\SecSet$ is a
\term[A pair of an Alice sector and a sector of the measured system; the joint sectors are ordered with the Alice sector slowest.]{joint sector};
the joint sectors are ordered with $t$ slowest, and
$\Hs{A\Qsys}=\bigoplus_{(t,s)}\Hs{A_t}\otimes\Hs{\Qsys_s}$, the $A$ index
first inside a block. The party partition labels $A$ with the party of the
factor tagged $A$, where there is one, and otherwise with the party label it
declares for $A$, as in the basic format; and it labels every tensor factor of
$\Hs{\Qsys_{f,s_f}}$ with the party of $f$.
\item \emph{(Alice's data and $\omega$ are block lists)} $\tau_A$, every
operator of Alice's measurement slot (DX4) and the right-hand side $\Theta_k$
of a general entry are lists of one array on $\Hs{A_t}$ per
$t\in\mathbb{T}$, and the list is the operator:
$\tau_A=\bigoplus_t\tau_A^{(t)}$, $M_x=\bigoplus_tM_x^{(t)}$,
$\Theta_k=\bigoplus_t\Theta_k^{(t)}$. A state $\omega$ on $A\Qsys$ is a list
of one array $\omega^{(t,s)}\in\Lin{A_t\Qsys_s}$ per joint sector,
$\omega=\bigoplus_{(t,s)}\omega^{(t,s)}$. Block diagonality along
$\mathbb{T}$ is a property of the type: no off-sector entry can be written.
By itself this typing licenses nothing; whether restricting $\omega$ to such
lists is sound is not settled by this definition.
\item \emph{(factorised family)} The \tref{outcome alphabet} is the product
$\Vset=\prod_{f\in\mathbb{F}}\Vset^{f}$ of the factors' ordered alphabets, in
the product order, the last factor fastest. What is declared is one family of
block lists $\{F^{f}_{v_f}\}_{v_f\in\Vset^{f}}$ per measured factor, against
that factor's own sector map, together with the rule
\begin{equation}\label{eq:productfamily}
  F_v\coloneqq\bigotimes_{f\in\mathbb{F}}F^{f}_{v_f},\qquad
  F^{(s)}_v=\bigotimes_{f\in\mathbb{F}}F^{f,(s_f)}_{v_f}
  \qquad\big(v=(v_f)_f,\ s=(s_f)_f\big).
\end{equation}
No entry of a product is declared. That $\{F_v\}_{v\in\Vset}$ is then a
POVM on $\Qsys$, block diagonal for the product sector map with the blocks of
\eqref{eq:productfamily}, follows for the exact numbers the factors declare:
a tensor product of Hermitian positive semidefinite operators is Hermitian
and positive semidefinite,
$\sum_{v}\bigotimes_fF^{f}_{v_f}=\bigotimes_f\sum_{v_f}F^{f}_{v_f}
=\idon{\Qsys}$ by multilinearity, and the tensor product distributes over the
direct sums $F^{f}_{v_f}=\bigoplus_{s_f}F^{f,(s_f)}_{v_f}$.
\item \emph{(the two forms of Alice's measurement slot)} In
\term[The form of Alice's measurement slot that holds the POVM itself.]{POVM form}
the slot holds $\{M_x\}_{x\in\Xset}$. In
\term[The trusted-source form of Alice's measurement slot that holds positive blocks summing to the marginal, and denotes a POVM-form declaration.]{ensemble form},
which exists in trusted-source mode only, it holds a family
$\{N_x\}_{x\in\Xset}$ of block lists along $\mathbb{T}$, every block
Hermitian and positive semidefinite, with
\begin{equation}\label{eq:ensembleform}
  \sum_{x\in\Xset}N^{(t)}_x=\tau^{(t)}_A
  \quad\text{and}\quad\tau^{(t)}_A\ne0
  \qquad\text{for every }t\in\mathbb{T},
\end{equation}
the sum formed entrywise in the \tref{declared order} of $\Xset$ and
rounding nowhere, and the second condition a comparison of declared entries with the
datum $0$. An ensemble-form declaration \emph{denotes} the POVM-form
declaration with the same remaining slots and
\begin{equation}\label{eq:ensemblepovm}
  M^{(t)}_x\coloneqq\big(\tau^{(t)}_A\big)^{-1/2}N^{(t)}_x
  \big(\tau^{(t)}_A\big)^{-1/2}
  +\frac{\tr N^{(t)}_x}{\tr\tau^{(t)}_A}\,\big(\idon{A_t}-P_t\big),
\end{equation}
where $P_t$ is the projector onto the support of $\tau^{(t)}_A$ and the
inverse square root is taken on that support and is zero off it. The
operators \eqref{eq:ensemblepovm} are defined and in general irrational; every statement made of an
ensemble-form declaration is made of the declaration it denotes. For a
POVM-form declaration in trusted-source mode put
$N_x\coloneqq\tau_A^{1/2}M_x\tau_A^{1/2}$; for the declaration an ensemble
form denotes, this returns the declared $N_x$.
\item \emph{(transcript protocol)} The announcement slot holds a
\term[The announcement slot of the general format: an ordered list of messages, each a function of its sender's readout and of the earlier messages.]{transcript protocol},
an ordered list of $L\ge1$ entries, each a
\term[One entry of a transcript protocol: a sender, a finite alphabet and a total table.]{message},
message $m$ carrying a sender $\pi_m\in\{\mathsf A,\mathsf B\}$, a finite
alphabet $\mathcal{C}_m$ and a total table
$a_m:\Oset{\pi_m}\times\mathcal{C}_1\times\dots\times\mathcal{C}_{m-1}
\to\mathcal{C}_m$. The joint announcement function is \emph{derived} from the
tables by composition,
\begin{equation}\label{eq:transcript}
  \kappa_m(x,v)\coloneqq a_m\big(\rd{\pi_m}(x,v),\,\kappa_1(x,v),\dots,
  \kappa_{m-1}(x,v)\big),\qquad
  a(x,v)\coloneqq\big(\kappa_m(x,v)\big)_{m=1}^{L}\in
  \Kann\coloneqq\prod_{m=1}^{L}\mathcal{C}_m ,
\end{equation}
and the acceptance-statistic map $\phi$, the sifting predicate $\varsigma$ of
\eqref{eq:chkey} and the holder's key map $r_{\hold}$ of D(viii) of
Definition~\ref{def:target} are maps on this $\Kann$. The two
maps $a_A$, $a_B$ of the basic format are the transcript protocol with $L=2$,
$\pi=(\mathsf A,\mathsf B)$, $\mathcal{C}_1=\CAset$, $\mathcal{C}_2=\CBset$,
$a_1=a_A$ and $a_2(o,\kappa_1)=a_B(o)$. A message is a function of its
sender's readout and of the earlier messages of the same round, and of
nothing else. No measurement of a round depends on a message of that round,
the measurements being the fixed families of (DX3) and (DX4); processing
across rounds is not a round message.
\item \emph{(party-scoped entries)} Each entry of $\CList$ carries the label
of one measured factor, its
\term[The measured factor that an entry of the constraint list is declared on.]{scope}
$f_k$, and its operator $\hat O_k$ is a block list against that factor's own
sector map, Hermitian block by block. Unless the entry is declared
\tref{party-local} it denotes the
\term[The operator on the whole measured system that a scoped entry's operator denotes: the identity on every other factor.]{lifted operator}
$\idon{}\otimes\hat O_k\otimes\idon{}$ on $\Qsys$, with $\hat O_k$ on the
tensor factor $\Qsys_{f_k}$, whose $s$-block is
$\idon{}\otimes\hat O_k^{(s_{f_k})}\otimes\idon{}$, and it is read in the forms stated
above for that lifted operator. With one measured factor the scope is that
factor and the lifted operator is $\hat O_k$ itself. An entry may instead be
declared party-local, in either source mode, carrying a real $c_k$; its
reading is the third form of Definition~\ref{def:homog}.
\end{enumerate}

\emph{What this module decides of the general format} is
Definition~\ref{def:exact} for the operator data and D(i)--D(viii) of
Definition~\ref{def:target} for the rest; that the family denoted by an
ensemble form is a POVM is Theorem~\ref{thm:sifting-and-key}\,(v).
\end{definition}
